\section{Contributions}

\paragraph{\textbf{Alessio Demo (Badge number: 2142885)}}
\begin{itemize}[leftmargin=*]
    \item \textbf{First Deadline:}  Developed function database connection and instance creation of DuckDB database manager for each dataset using the ingest files (.duckdb files).
    Developed function for retrieving the SQL queries from .sql files, processing the queries and executing them on DuckDB database instance and finally store it in a json file.
    \item \textbf{Second Deadline:} Developed the Baseline NL function for interacting with the LLM. "Early-stage" Palimpzest baseline implementation.
\end{itemize}

\paragraph{\textbf{Francesco Pivotto (Badge number: 2158296)}}
\begin{itemize}[leftmargin=*]
    \item \textbf{First Deadline:} Implemented a modular project structure by relocating core logic to \texttt{src/} and consolidating configuration and database management logic within dedicated utility modules (\texttt{config/loaders.py}). 
    Established the framework for multi-provider support by adding initial integration for the \texttt{Grok LLM}, including dedicated configuration points (\texttt{.env}, \texttt{config.yaml}, \texttt{loaders.py}).
    
    \item \textbf{Second Deadline:} Completed initial implementation of the SQL baseline, including the core logic for connecting to \texttt{DuckDB} and structuring query execution via the Watsonx provider.
    Fixed critical issues in environment variable loading, ensuring robust and standardized retrieval of secrets and settings via the central configuration object (using \texttt{loaders.py}).
    Modularized the \texttt{sql\_baseline} module, incorporating the LCEL Pattern to prepare for dynamic LLM selection.
    Centralized general utility functions by moving helpers into \texttt{baseline\_tools}.
    Introduced the dedicated \textbf{\texttt{llm\_wrapper}} class to handle all LLM connection logic, API parameter management, and calls.
\end{itemize}

\paragraph{\textbf{Giorgia Amato (Badge number: ...)}}